\subsection{The approximate solution}\label{cc:uf:sec:approximate}
Fix an integer $P\ge2$. The profiles $W,\mathsf P$ of
Proposition~\ref{cc:uf:hier:prop:identity} at depth $P$ give the
Eulerian velocity and pressure
\[
 V=U+W^\natural,\qquad \Pi=p_U+\mathsf P^\natural,\qquad
 W=W_{\rm osc}+\bar W,\quad \bar W=\sum_{p=2}^PB_p .
\]
Fix an absolute integer $b\ge1$, at least as large
as the exponents $C$ of Theorem~\ref{cc:uf:hier:thm:depth} and of
Lemma~\ref{cc:uf:hier:lem:physical}, and sufficiently large for the
finitely many estimates of this subsection and the next. Put
\[
 C_{\rm step}=\Gamma_P\mathsf D^b,
\]
where $\Gamma_P\ge1$ depends only on $P$ and the fixed profiles.
We take it sufficiently large compared with every $C_{2,P}$ in this
subsection and the next. The estimates of the residual, the pressure
Hessian and the flow use the profile bounds through order two. The
truncation uses the velocity bound in $\Holder{P}$ and the tail estimate
\eqref{cc:uf:trunc:eq:tail}. From here on we assume
\begin{equation}\label{cc:uf:step:eq:small}
 C_{\rm step}\lambda^{-1/12}\le1 .
\end{equation}
Since $C_{\rm step}\ge\mathsf D^C$ for the exponent $C$ of
Theorem~\ref{cc:uf:hier:thm:depth}, the hypothesis of that theorem
holds.

To estimate the velocity and residual in Eulerian coordinates, we evaluate
each profile at $a=X(t,\cdot)^{-1}(x)$ and $s=k_*\cdot a$.

\begin{lemma}\label{cc:uf:hier:lem:physical}
Let \(g_0(t,a,s)\) be smooth and independent of the phase at the labels
outside \(\supp g\). Assume its phase average extends smoothly and
periodically in $a$ across the cell faces. There is an absolute exponent
$C$ such that, for every integer \(m\ge0\) and every $t\in[t_0,T]$,
\begin{align}
 \sum_{|\beta|=m}\|\partial_x^\beta g_0^\natural(t)\|_{L^\infty}
 &\le C_m(\mathsf D^C\lambda^{3/2})^m\gamma(t)
                  \|g_0\|_{[m],\gamma},
 \label{cc:uf:hier:eq:physical-infty}\\
 \|g_0^\natural(t)\|_{\Holder{m}}
 &\le C_m(\mathsf D^C\lambda^{3/2})^{m+1}\gamma(t)
                  \|g_0\|_{[m+1],\gamma}.
 \label{cc:uf:hier:eq:physical}
\end{align}
Both bounds also hold on $[0,T]$ with \(\gamma\) replaced by one, and for
a phase-independent field with \(\mathsf D^C\lambda^{3/2}\) replaced by
\(\mathsf D^C\).
\end{lemma}
\begin{proof}
The field \(\widetilde{\mathbf d}_ig_0\) is again independent of the
phase outside \(\supp g\), which lies in the interior of the cell, so
\(g_0^\natural\) and its derivatives extend smoothly and periodically.
The spatial chain rule in \eqref{cc:uf:hier:eq:graph} and
\(\widetilde{\mathbf d}=(D_aX)^{-T}(\nabla_a+k_*\partial_s)\) give
\[
 \sum_i\|\widetilde{\mathbf d}_ig_0\|_{[n],\gamma}
 \le C_n\mathsf D^C\lambda^{3/2}\|g_0\|_{[n+1],\gamma},\qquad n\ge0,
\]
by Lemma~\ref{cc:uf:trunc:lem:fields} and
\eqref{cc:uf:hier:eq:coefficient-bound}: a label derivative costs
$\ell_*^{-1}$, a fixed power of $\mathsf D$, and $k_*\partial_s$ costs
$|k_*|h^{-1}\le\mathsf D\lambda^{3/2}$, by
\eqref{cc:uf:tower:eq:wave-vector}.
Iterate the displayed bound; at each time the supremum of a field is at
most $\gamma(t)$ times its weighted norm of order zero. This proves
\eqref{cc:uf:hier:eq:physical-infty}. The bound
\(\|v\|_{\Holder{0}}\le C(\|v\|_\infty+\|\nabla v\|_\infty)\),
applied to derivatives through order \(m\), proves
\eqref{cc:uf:hier:eq:physical} with one additional derivative. A
phase-independent field has
\(\widetilde{\mathbf d}g_0=\mathbf dg_0\), proving the last assertion.
\end{proof}

In Section~\ref{cc:uf:sec:tower} the oscillations are summed in the
supremum norm, while their higher derivatives control the Fourier tail
in \eqref{cc:uf:tower:eq:tail-small}. The means can extend across the
torus and back to time zero, so their sum must converge at every fixed
derivative order.

\begin{lemma}\label{cc:uf:step:lem:approximate}
The field \(V\) is smooth, odd and
divergence free, \(\Pi\) is smooth, and
\begin{equation}\label{cc:uf:step:eq:approximate}
 V_t+V\cdot\nabla V+\nabla\Pi=R_V,\qquad
 R_V=R_P^\natural+R_{\widetilde\chi}^\natural .
\end{equation}
The oscillation satisfies
\begin{equation}\label{cc:uf:step:eq:physical-sup}
 \|W_{\rm osc}^\natural\|_{C_tL^\infty}
 \le C_{0,P}\mathsf D^b\lambda^{-5/4},
\end{equation}
and for every \(m\ge0\),
\begin{equation}\label{cc:uf:step:eq:physical}
\begin{gathered}
 \|W_{\rm osc}^\natural\|_{C_t\Holder{m}}
 \le C_{m,P}\mathsf D^{b(m+2)}\lambda^{3(m+1)/2-5/4},\qquad
 \|\bar W^\natural\|_{C_t\Holder{m}}
 \le C_{m,P}\mathsf D^{b(m+2)}\lambda^{-5/2},\\
 \|R_V\|_{C_t\Holder{1}}
 \le C_{\rm step}\lambda^2
   \bigl(\lambda^{-2P/3}+\|\widetilde\chi'\gamma\|_\infty\bigr),
\end{gathered}
\end{equation}
where the last supremum is over $[t_0,T]$.
\end{lemma}
\begin{proof}
By \eqref{cc:uf:hier:eq:graph}, at \(x=X(t,a)\) one has
\((\partial_t+U\cdot\nabla)W^\natural=(W_t)^\natural\),
\(W^\natural\cdot\nabla U=(MW)^\natural\),
\(W^\natural\cdot\nabla W^\natural=((W\cdot\widetilde{\mathbf d})W)^\natural\) and
\(\nabla\mathsf P^\natural=(\widetilde{\mathbf d}\mathsf P)^\natural\). Adding
the definition of \(R_U\) and \eqref{cc:uf:hier:eq:residual} gives
\eqref{cc:uf:step:eq:approximate}, because \((R_U\circ X)^\natural=R_U\):
the residual of the parent is canceled exactly. Also
\(\widetilde{\mathbf d}\cdot W=0\) gives \(\diver V=0\).
Since \(\supp g\) lies in the interior of the cell, the oscillatory
terms in \(W^\natural\) and \(\mathsf P^\natural\) extend smoothly and
periodically for every \(k_*\). Proposition~\ref{cc:uf:hier:prop:identity}
gives the smoothness of \(V\) and \(\Pi\), and \(W^\natural\) is odd by
the symmetry of the construction of Subsection~\ref{cc:uf:sec:hierarchy}
under \((a,s)\mapsto(-a,-s)\).

Since $\gamma\le1$, the supremum of
$W_{\rm osc}^\natural$ is at most $\|W_{\rm osc}\|_{[0],\gamma}$, and
\eqref{cc:uf:hier:eq:sum} gives \eqref{cc:uf:step:eq:physical-sup}.
The bound \eqref{cc:uf:hier:eq:physical} of
Lemma~\ref{cc:uf:hier:lem:physical}, applied to $W_{\rm osc}$ and
$\bar W$ with \eqref{cc:uf:hier:eq:sum} at the order $m+1$, proves the
first two bounds of \eqref{cc:uf:step:eq:physical}.
By Proposition~\ref{cc:uf:hier:prop:identity}, $R_P$ is independent of
the phase outside $\supp g$, so \eqref{cc:uf:hier:eq:physical} at $m=1$,
with $\gamma$ replaced by one, and \eqref{cc:uf:hier:eq:tail-est} bound
$R_P^\natural$ in $C_t\Holder{1}$ by
$C_{2,P}\mathsf D^C\lambda^{7/4}\lambda^{-2P/3}$. By
\eqref{cc:uf:hier:eq:size} and \eqref{cc:uf:hier:eq:completion-bound},
$\|A_1+C[A_1]\|_{[2],\gamma}\le C\mathsf D^C\lambda^{-5/4}$. At each time
$t$, $R_{\widetilde\chi}$ is $\widetilde\chi'(t)$ times this field, so
\eqref{cc:uf:hier:eq:physical} at that time bounds
$R_{\widetilde\chi}^\natural(t)$ in $\Holder{1}$ by
$C\mathsf D^C\lambda^{7/4}|\widetilde\chi'(t)|\gamma(t)$. This proves
the last bound of \eqref{cc:uf:step:eq:physical}.
\end{proof}

In \eqref{cc:uf:step:eq:physical} the depth $P$ controls the expansion
remainder. For a later wave, \eqref{cc:uf:tower:eq:switch-small}
bounds $\gamma$ on the interval where the time cutoff varies.

The next mean equation requires an upper bound on the added pressure
Hessian, since the pressure Hessian enters the displacement form
\eqref{cc:uf:mean:jacobi-id} with a minus sign. The next wave also needs
the velocity gradient at the origin. Once $\widetilde\chi=1$, the
rank-one term in \eqref{cc:uf:tower:eq:central-return} supplies the
wave contribution to the gradient in \eqref{cc:uf:tower:eq:split}.

\Needspace{8\baselineskip}
\begin{lemma}\label{cc:uf:step:lem:pressure-center}
On \([0,T]\times\mathbb T^3\),
\begin{equation}\label{cc:uf:step:eq:hessian}
 D_x^2\mathsf P^\natural\le C_{\rm step}\lambda^{-1/4}\Id .
\end{equation}
At \(x=0\), for \(t_0\le t\le T\),
\begin{equation}\label{cc:uf:tower:eq:central-return}
 \nabla V(t,0)=\nabla U(t,0)
   +\widetilde\chi(t)\,\mathbf a(t,0)\otimes\kappa(t,0)+E_c(t),
 \qquad
 |E_c|\le\lambda^{-1/3} .
\end{equation}
\end{lemma}
\begin{proof}
Write the leading pressure as \(\mathsf p_1=-c_{\mathbf a}P_h\), with
\(c_{\mathbf a}=2\widetilde\chi g(\kappa\cdot M\mathbf a)/|\kappa|^2\). Since
$D_x^2\mathsf p_1^\natural
=(\widetilde{\mathbf d}\,\widetilde{\mathbf d}\,\mathsf p_1)^\natural$ and
\(\mathbf d_i\kappa_j=\mathbf d_j\kappa_i\), with \(a=X(t,\cdot)^{-1}(x)\)
and \(s=k_*\cdot a\) on the right,
\begin{equation}\label{cc:uf:exact:eq:primary-hess}
 D_x^2\mathsf p_1^\natural=
 -c_{\mathbf a}F_h'\kappa\otimes\kappa
 +(\mathbf d\partial_s\mathsf p_1)\otimes\kappa+
 \kappa\otimes(\mathbf d\partial_s\mathsf p_1)
 +(\partial_s\mathsf p_1)\mathbf d\kappa+\mathbf d^2\mathsf p_1 .
\end{equation}
The product $\widetilde\chi g$ takes values in $[0,1]$ and vanishes off
$[t_0,T]\times\supp g$, where $\kappa\cdot M\mathbf a>0$ by
\eqref{cc:uf:tower:eq:near-origin}. Hence
\(0\le c_{\mathbf a}|\kappa|^2\le2|\kappa||M||\mathbf a|\). With
\(F_h'\ge-Ch\), the first tensor on the right is therefore at most
\(Chc_{\mathbf a}|\kappa|^2\Id\). On $[t_0,T]\times\supp g$,
\eqref{cc:uf:tower:eq:near-origin} gives
$|\mathbf a|\le2\lambda^{-3/4}$, while $|\kappa|\le\mathsf D^2\lambda$
and $|M|\le\mathsf D$ by \eqref{cc:uf:tower:eq:wave-vector} and
\eqref{cc:uf:tower:eq:parent-bound}. So
\(c_{\mathbf a}|\kappa|^2\le4\mathsf D^3\lambda^{1/4}\) and, since
$h=\lambda^{-1/2}$, the first tensor is at most
\(C\mathsf D^3\lambda^{-1/4}\Id\). By
\eqref{cc:uf:trunc:eq:product}, \eqref{cc:uf:trunc:eq:deriv} and
\eqref{cc:uf:hier:eq:coefficient-bound}, the sum of the remaining
tensors is at most
$C\mathsf D^C\lambda^{3/2}\|\mathsf p_1\|_{[2]}$ in supremum norm, and
\(\lambda\|\mathsf p_1\|_{[2]}\le C\mathsf D^C\lambda^{-5/4}\)
by \eqref{cc:uf:hier:eq:size}. So their sum is at most
$C\mathsf D^C\lambda^{-3/4}$.

For the higher oscillatory pressures, \eqref{cc:uf:hier:eq:Vbound}
gives the least decay in $\lambda$ at $p=2$, where it bounds
$\mathsf p_2$ by $C_{2,P}\mathsf D^C\lambda^{-1-23/12}$.
Two Eulerian derivatives in \eqref{cc:uf:hier:eq:physical-infty}
would cost $\lambda^3$, leaving $\lambda^{1/12}$. To bound the first
derivative, use the transverse equation \eqref{cc:uf:hier:eq:pressure-from-trans}.
For $2\le p\le P$ it gives
\[
 \nabla_x\mathsf p_p^\natural
   =(\mathbf d\mathsf p_p+\kappa\partial_s\mathsf p_p)^\natural,
 \qquad
 \kappa\partial_s\mathsf p_p=P_\kappa(f_p+f_p^S-2MV_p).
\]
The label derivative costs a power of $\mathsf D$, and the phase term
is bounded by the source and velocity in \eqref{cc:uf:hier:eq:Vbound}.
The coefficient estimates therefore give
\[
 \Big\|\widetilde{\mathbf d}\sum_{p=2}^P\mathsf p_p\Big\|_{[1],\gamma}
 \le C_{2,P}\mathsf D^C\lambda^{-23/12}.
\]
The oscillatory pressures are supported in $\supp g$ in Lagrangian coordinates.
One further Eulerian derivative costs $\lambda^{3/2}$ up to a power
of $\mathsf D$, so \eqref{cc:uf:hier:eq:physical-infty} gives
\begin{equation}\label{cc:uf:exact:eq:higher-hess}
 \Big\|D_x^2\sum_{p=2}^P\mathsf p_p^\natural\Big\|_\infty
 \le C_{2,P}\mathsf D^C\lambda^{-5/12}.
\end{equation}
For the mean pressures, add the mean equations
\eqref{cc:uf:hier:eq:mean} of the levels $2\le p\le P$. In Eulerian coordinates the
sum is the linearized Euler equation \eqref{cc:uf:mean:physical} for
$\bar W^\natural$ over $U$, with the pressure
$\sum_p\bar q_p^\natural$ and the source $(\sum_p\bar f_p)^\natural$,
and its divergence is
\[
 \Delta\sum_{p=2}^P\bar q_p^\natural
 =\diver\Bigl(\sum_{p=2}^P\bar f_p\Bigr)^\natural
  -2\tr\bigl(\nabla U\,\nabla\bar W^\natural\bigr).
\]
The operator $D^2\Delta^{-1}$ is bounded on $\Holder{0}$. By
\eqref{cc:uf:hier:eq:physical} for phase-independent fields and
\eqref{cc:uf:hier:eq:sum}, the source and $\bar W^\natural$ are at most
$C_{2,P}\mathsf D^C\lambda^{-5/2}$ in $C_t\Holder{1}$, and
$\|\nabla U\|_{C_t\Holder{0}}\le C\mathsf D\norm{\nabla U}_\infty
\le C\mathsf D^2$ by \eqref{cc:uf:trunc:eq:lp} and
\eqref{cc:uf:tower:eq:parent-bound}. Hence the Hessian of the
sum of the mean pressures is at most
$C_{2,P}\mathsf D^C\lambda^{-5/2}$ in supremum norm. Together with the
bounds for the tensors of \eqref{cc:uf:exact:eq:primary-hess} and with
\eqref{cc:uf:exact:eq:higher-hess}, this proves
\eqref{cc:uf:step:eq:hessian}.

To compute the velocity gradient at $x=0$, evaluate at $a=0$ and $s=0$,
since the flow of $U$ fixes the origin. As
\(g(0)=F_h'(0)=1\) and \(F_h(0)=0\), the chain rule
\eqref{cc:uf:hier:eq:graph} gives
\(D_xV_1^\natural(t,0)=\widetilde\chi(t)\mathbf a(t,0)\otimes\kappa(t,0)\).
Hence
\[
 E_c(t)=D_x(W_{\rm osc}-V_1)^\natural(t,0)+D_x\bar W^\natural(t,0).
\]
The field \(W_{\rm osc}-V_1=\sum_{p=2}^PV_p+\sum_{p=2}^{P+1}C_p\) is supported in
\(\supp g\). By \eqref{cc:uf:hier:eq:Vbound} and
\eqref{cc:uf:hier:eq:Cbound}, its norm \(\|\cdot\|_{[1],\gamma}\) is at
most \(C_{2,P}\mathsf D^C\lambda^{-23/12}\), the bound for $V_2$,
and \(\|\bar W\|_{[1]}\le C_{2,P}\mathsf D^C\lambda^{-5/2}\) by
\eqref{cc:uf:hier:eq:sum}. The bound \eqref{cc:uf:hier:eq:physical-infty}
at \(m=1\) therefore gives
\(|E_c|\le C_{2,P}\mathsf D^C\lambda^{-5/12}
=C_{2,P}\mathsf D^C\lambda^{-1/12}\cdot\lambda^{-1/3}\). Since
$b$ is at least the exponent $C$ of this bound and
$\Gamma_P\ge C_{2,P}$, \eqref{cc:uf:step:eq:small} makes it at
most $\lambda^{-1/3}$.
This proves \eqref{cc:uf:tower:eq:central-return}.
\end{proof}

To allow for the difference introduced by frequency truncation, let $Y$
be the flow of $V+v$ for a small divergence-free field $v$. A direct Gronwall
estimate for $D_aY$ puts the large velocity gradient in an exponential.
Write $Y(t,a)=X(t,\mathcal A_t(a))$ and let the label and phase evolve
as independent variables. Incompressibility expresses the oscillatory
phase velocity through a label divergence, giving a small velocity
gradient in the scaled variables. We then differentiate the change of
label on the graph $s=k_*\cdot a$ to recover $D\mathcal A_t$.

\begin{lemma}\label{cc:uf:step:lem:flow-bounds}
Let $v$ be a smooth divergence-free field on $[0,T]\times\mathbb T^3$
with $\norm v_{C_tC^2}\le\lambda^{-5/2}$. The flow $Y$ of $V+v$ from
time zero satisfies
\begin{equation}\label{cc:uf:step:eq:deformation}
 \sup|(D_aY)^{\pm1}|\le C_{\rm step}\lambda^{1/4}.
\end{equation}
\end{lemma}
\begin{proof}
The field $v\circ X$ is independent of the phase, and
$\mathbf d\cdot(v\circ X)=(\diver v)\circ X=0$, so
$\widetilde{\mathbf d}\cdot(W+v\circ X)=0$. By the chain rule and
\eqref{cc:uf:hier:eq:coefficient-bound} for $D_aX$,
$\|v\circ X\|_{[2]}\le C\mathsf D^C\norm v_{C_tC^2}$.
Put $c=(D_aX)^{-1}(W+v\circ X)$ and write $c=c_{\rm osc}+\bar c$ for its
oscillatory and phase-independent parts, with $c_{\rm osc}$ repeated over
the lattice by $(a,s)\mapsto(a+n,s+k_*\cdot n)$,
$n\in\mathbb Z^3$. On the graph $s=k_*\cdot a$, the field $c$ is then
$(D_aX)^{-1}((W^\natural+v)\circ X)$. The Piola identity turns the
solenoidality of $W+v\circ X$ into $(\nabla_a+k_*\partial_s)\cdot c=0$,
and the zero-mean phase primitive of its oscillatory part gives
$k_*\cdot c_{\rm osc}=-\partial_s^{-1}\diver_a c_{\rm osc}$. With
\eqref{cc:uf:hier:eq:sum} and \eqref{cc:uf:hier:eq:coefficient-bound},
the norms without weight therefore satisfy
\[
 \|c_{\rm osc}\|_{[2]}\le C_{2,P}\mathsf D^C\lambda^{-5/4},\qquad
 \|\bar c\|_{[2]}\le C_{2,P}\mathsf D^C\lambda^{-5/2},\qquad
 \|k_*\cdot c\|_{[1]}\le C_{2,P}\mathsf D^C\lambda^{-5/4};
\]
in the last bound the oscillatory part costs one label derivative and
the mean part the factor $|k_*|\le\mathsf D\lambda$ of
\eqref{cc:uf:tower:eq:wave-vector}.

Let \(\mathscr L_t\), with first component \(a+\eta_a(t,a,s)\), be the
flow from time zero of
\begin{equation}\label{cc:uf:exact:eq:particle}
 \dot a=c(t,a,s),\qquad \dot s=k_*\cdot c(t,a,s).
\end{equation}
In the coordinates $(a/\ell_*,s/h)$ this vector field is
$(c/\ell_*,k_*\cdot c/h)$, and its gradient is at most
$C(\ell_*^{-1}\|c\|_{[1]}+h^{-1}\|k_*\cdot c\|_{[1]})
\le C_{2,P}\mathsf D^C\lambda^{-3/4}
\le C_{2,P}\mathsf D^C\lambda^{-1/12}$. Since $b$ is at least the
exponent $C$ of this bound, \eqref{cc:uf:step:eq:small} makes it at
most $C_{2,P}/\Gamma_P$, a small absolute constant since $\Gamma_P$ is
large.
Since $t\le T\le1$, Gronwall's inequality then bounds the
gradients of $\mathscr L_t$ and of $\mathscr L_t^{-1}$ in these
coordinates by $2$. The scaled gradient of $c$ is at most
$\|c\|_{[1]}$, so integrating $\dot\eta_a=c$ along the flow gives
\[
 \ell_*|\partial_a\eta_a|+h|\partial_s\eta_a|
 \le C\|c\|_{[1]}\le C_{2,P}\mathsf D^C\lambda^{-5/4},
\]
and the displacement of $\mathscr L_t^{-1}$ satisfies the same bound.

Differentiating $Y=X(t,\mathcal A)$ gives
$\partial_t\mathcal A=((D_aX)^{-1}((W^\natural+v)\circ X))(t,\mathcal A)$.
Since $s-k_*\cdot a$ is preserved by \eqref{cc:uf:exact:eq:particle},
the graph $s=k_*\cdot a$ is invariant, so
$\mathcal A_t(a)=a+\eta_a(t,a,k_*\cdot a)$, the first
component of \(\mathscr L_t\) on this graph. The same formula with the
displacement of \(\mathscr L_t^{-1}\) gives the inverse map.
Thus $D\mathcal A_t-\Id=\partial_a\eta_a+\partial_s\eta_a\otimes k_*$,
and with the corresponding identity for the inverse,
\[
 |D\mathcal A_t^{\pm1}-\Id|
 \le\bigl(\ell_*^{-1}+|k_*|h^{-1}\bigr)C_{2,P}\mathsf D^C\lambda^{-5/4}
 \le C_{2,P}\mathsf D^C\lambda^{1/4}.
\]
Then \(D_aY=(D_aX\circ\mathcal A)D\mathcal A\) and
\((D_aY)^{-1}=(D\mathcal A)^{-1}((D_aX)^{-1}\circ\mathcal A)\), with
the bound \eqref{cc:uf:tower:eq:parent-bound} for the deformation of
the parent, give \eqref{cc:uf:step:eq:deformation}.
\end{proof}

\Needspace{11\baselineskip}
\subsection{Frequency truncation}\label{cc:uf:trunc:sec:truncation}
We truncate $V$ with the smooth frequency cutoff $S_\Lambda$ of
Appendix~\ref{cc:app:cutoff}, as in the mollification step
of~\cite[Section~2.4]{BDLSVOnsager}. The cutoff bounds all higher
derivatives of the next parent in terms of $\Lambda$ and its supremum.
We choose its pressure $p$ so that its residual is divergence free.
The equation for the expansion pressure $\Pi$ contains $\diver R_V$
by \eqref{cc:uf:step:eq:approximate}, so the bound for $D^2(p-\Pi)$
includes both the residual of $V$ and the Fourier tail. The iteration
uses the gradient bound on the new residual and the upper bound on
$D^2(p-p_U)$ to propagate \eqref{cc:uf:tower:eq:parent-coercive}.

\begin{proposition}\label{cc:uf:step:thm:step}
Let $\Lambda\ge1$, and with the depth $P$ of the expansion put
\begin{equation}\label{cc:uf:step:eq:truncation}
 u=S_\Lambda V,\qquad v_\Lambda=u-V,\qquad
 \varepsilon_\Lambda=\Lambda^{2-P}
 \bigl(1+\norm V_{C_t\Holder{P}}\bigr)^2 .
\end{equation}
Let $p$ be the function of zero mean with
$\Delta p=-\tr((\nabla u)^2)$, and let
$R_u=u_t+u\cdot\nabla u+\nabla p$. Then $u$ is smooth, odd, divergence
free and supported on frequencies at most $2\Lambda$, its residual $R_u$ is odd,
divergence free and supported on frequencies at most $4\Lambda$, and
\begin{equation}\label{cc:uf:step:eq:truncated}
 \norm{v_\Lambda}_{C_t\Holder{2}}\le C_{2,P}\varepsilon_\Lambda,\qquad
 \norm{R_u}_{C_t\Holder{1}}+\norm{D^2(p-\Pi)}_{C_t\Holder{0}}
 \le C\norm{R_V}_{C_t\Holder{1}}+C_{2,P}\varepsilon_\Lambda .
\end{equation}
Moreover, on $[0,T]\times\mathbb T^3$,
\begin{equation}\label{cc:uf:step:eq:pressure-error}
 D^2(p-p_U)\le\bigl(C_{\rm step}\lambda^{-1/4}
   +C\norm{R_V}_{C_t\Holder{1}}+C_{2,P}\varepsilon_\Lambda\bigr)\Id .
\end{equation}
\end{proposition}
\begin{proof}
By Lemma~\ref{cc:uf:trunc:lem:cutoff}, $u$ and $v_\Lambda$ are odd and
divergence free and $u$ is supported on frequencies at most $2\Lambda$.
Since $\diver(u\cdot\nabla u)=\tr((\nabla u)^2)$, the residual is
$R_u=\partial_tu+\PP\diver(u\otimes u)$; it is odd, divergence free and
supported on frequencies at most $4\Lambda$.
Applying $S_\Lambda$ and then $\PP$ to \eqref{cc:uf:step:eq:approximate}
gives
$\partial_tu+\PP\diver S_\Lambda(V\otimes V)=\PP S_\Lambda R_V$.
Subtracting this from the formula for $R_u$ gives
\begin{equation}\label{cc:uf:trunc:eq:residual}
 R_u=\PP\bigl[S_\Lambda R_V
   +\diver\bigl(v_\Lambda\otimes u+V\otimes v_\Lambda
                +(\Id-S_\Lambda)(V\otimes V)\bigr)\bigr].
\end{equation}
To bound $R_u$ in $\Holder{1}$ we need the three quadratic tensors in
\eqref{cc:uf:trunc:eq:residual} in $\Holder{2}$. This accounts for
the power $\Lambda^{2-P}$ and the square in $\varepsilon_\Lambda$.
The tail bound \eqref{cc:uf:trunc:eq:tail} at the order $P$ gives
$\norm{v_\Lambda}_{\Holder{s}}
\le C_{2,P}\Lambda^{s-P}\norm V_{\Holder{P}}$ for
$s\le2$, and in particular the first bound of
\eqref{cc:uf:step:eq:truncated}. By Leibniz' rule and the algebra
inequality for the norm of $\Holder{0}$, for every integer $m\ge0$,
\begin{equation}\label{cc:holder:product-growth}
 \norm{fg}_{\Holder{m}}\le C_m\norm f_{\Holder{m}}\norm g_{\Holder{m}}.
\end{equation}
Hence the three tensors on the right of
\eqref{cc:uf:trunc:eq:residual} are at most
$C_{2,P}\Lambda^{2-P}\norm V_{\Holder{P}}^2\le C_{2,P}\varepsilon_\Lambda$ in
$\Holder{2}$. The periodic
multipliers $\PP$ and $D^2\Delta^{-1}$ of order zero are bounded
on $\Holder{0}$ and commute with derivatives, as does $S_\Lambda$; this
bounds $R_u$ in $\Holder{1}$.
For the pressure, the divergence of \eqref{cc:uf:step:eq:approximate} is
$\Delta\Pi=\diver R_V-\tr((\nabla V)^2)$. Hence
\[
 \Delta(p-\Pi)
 =-\tr(\nabla v_\Lambda\nabla V)-\tr(\nabla u\nabla v_\Lambda)
  -\diver R_V,
\]
and $\norm{D^2(p-\Pi)}_{\Holder{0}}
\le C\norm{v_\Lambda}_{\Holder{1}}(\norm V_{\Holder{1}}+\norm u_{\Holder{1}})
+C\norm{R_V}_{\Holder{1}}$. This proves \eqref{cc:uf:step:eq:truncated}.
Since $p-p_U=(p-\Pi)+\mathsf P^\natural$,
Lemma~\ref{cc:uf:step:lem:pressure-center} then gives
\eqref{cc:uf:step:eq:pressure-error}.
\end{proof}
